% Part: lambda-calculus
% Chapter: lambda-definability
% Section: truth-values

\documentclass[../../../include/open-logic-section]{subfiles}

\begin{document}

\olfileid{lam}{ldf}{tvr}
\olsection{真值与关系}

在纯λ演算中，我们可以如下编码真值：
\begin{align*}
  \fn{true} & \ident \lambd[x][\lambd[y][x]]\\
  \fn{false} & \ident \lambd[x][\lambd[y][y]]
\end{align*}

真值表示为\emph{选择器}，即接受两个自变元并返回其中之一的函数。真值
$\fn{true}$ 选择其第一个自变元，而 $\fn{false}$ 选择其第二个自变元。
例如，$\fn{true}\, M N$ 总是归约为 $M$，而 $\fn{false}\, M N$ 总是归约为~$N$。

\begin{defn}
称关系 $R \subseteq \Nat^n$ 是 $\lambd$\emph{-可定义的}，如果存在项~$R$ 使得
\begin{align*}
  R\, \num{n_1} \dots \num{n_k} & \bred \fn{true}
  \intertext{若 $R(n_1, \dots, n_k)$；否则}
  R\, \num{n_1} \dots \num{n_k} & \bred \fn{false}
\end{align*}。
\end{defn}

例如，仅对 $0$ 成立的关系 $\fn{IsZero} = \{0\}$，可由下式给出 $\lambd$-定义：
\[
  \fn{IsZero} \ident \lambd[n][n (\lambd[x][\fn{false}])\, \fn{true}].
\]
它的工作原理是什么？由于 邱奇 数被定义为迭代器（即将其第一个自变元应用到第二个自变元上 $n$~次的函数），
我们将初始值设为 $\fn{true}$，并且在迭代的每一步，无论上一次迭代结果如何，
都返回 $\fn{false}$。此步骤将对初始值应用 $n$ 次，而结果为
$\fn{true}$ 当且仅当该步骤根本没有被应用，即 $n = 0$ 时。

基于这种真值表示，我们可以进一步定义一些真值函数。下面给出两个，即否定与合取的表示：
\begin{align*}
  \fn{Not} & \ident \lambd[x][x\, \fn{false}\, \fn{true}]\\
  \fn{And} & \ident \lambd[x][\lambd[y][xy \,\fn{false}]]
\end{align*}
函数“$\fn{Not}$”接受一个自变元，若该自变元是 $\fn{false}$ 则返回 $\fn{true}$，
若该自变元是~$\fn{true}$ 则返回~$\fn{false}$。函数“$\fn{And}$”接受两个
真值作为自变元，且当且仅当两个自变元均为~$\fn{true}$ 时才返回 $\fn{true}$。
真值表示为选择器（如上所述），因此当 $x$ 是真值并被应用于两个自变元时，
若 $x$ 为 $\fn{true}$，结果为第一个自变元，否则为第二个自变元。
现在，$\fn{And}$ 接受其两个自变元 $x$ 和 $y$，然后将 $y$ 和 $\fn{false}$
传递给其第一个自变元~$x$。假设 $x$ 是真值，那么若 $x$
是 $\fn{true}$，结果将求值为~$y$；若 $x$
是 $\fn{false}$，结果将求值为 $\fn{false}$，
而这正是我们所期望的。

注意，我们这里假设仅将真值用作 $\fn{And}$ 的自变元。若传入其他项，
则结果（即范式，若存在的话）很可能不是真值。

\begin{prob}
  利用将真值编码为 $\lambd$ 项的方法，定义表示相容析取与不相容析取的真值函数
  $\fn{Or}$ 和 $\fn{Xor}$。
\end{prob}

\end{document}